\mode*
\section{Introduction}

The Unix shell is one of the first tools computing students meet: it is how
they access the university's systems, hand in their code, and later drive
the tools of their trade.
It is widely taught---and yet, in contrast to introductory programming,
there is little research on how students learn the shell or on how to teach
it \autocite{Winder2024ShellTutor}.
What we do know suggests that the difficulties are real: producing composite
shell commands is demanding even for users with some experience, and what
helps novices (concrete information) differs qualitatively from what helps
experts (abstract information) \autocite{Doane1992UnixPrompts}.
Moreover, users in general persistently fail to adopt faster expert methods
when a familiar interface is available \autocite{Cockburn2014NoviceExpert},
so students will not move from the graphical interface to the shell merely
because the shell is more powerful.

\mode<presentation>{%
\begin{frame}
  \begin{block}{The shell}
    \begin{itemize}
      \item Widely taught, little researched
        \autocite{Winder2024ShellTutor}.
      \item Composite commands are hard; novices need \emph{concrete}
        information \autocite{Doane1992UnixPrompts}.
      \item Users don't switch interfaces voluntarily
        \autocite{Cockburn2014NoviceExpert}.
    \end{itemize}
  \end{block}
\end{frame}%
}

In this work we apply the approach of our companion paper on introductory
programming \autocite{VTProgMisconceptions}: we treat documented
misconceptions and difficulties as the phenomenographic observations for a
variation-theory analysis \autocite{NCOL}, identify the critical aspects
that learners must discern, and outline patterns of variation to teach them.
We cover the shell's command model, the hierarchical file system,
composition (pipes, redirection, scripting), and the remote shell (SSH).

We ask the following research questions.

\begin{frame}
\begin{restatable}{question}{rqmisconceptions}\label{rq:misconceptions}
  Which misconceptions and difficulties with the Unix shell---its command
  model, the file system, composition, and remote use---does the literature
  document?
\end{restatable}

\begin{restatable}{question}{rqaspects}\label{rq:aspects}
  Which critical aspects of using the shell do these misconceptions and
  difficulties point to?
\end{restatable}

\begin{restatable}{question}{rqpatterns}\label{rq:patterns}
  Which patterns of variation can make these critical aspects discernible
  to learners?
\end{restatable}
\end{frame}

Our contributions are:
a review of the literature on learners' difficulties with the shell,
answering \cref{rq:misconceptions};
a variation-theory analysis identifying the critical aspects,
answering \cref{rq:aspects};
and a set of patterns of variation for teaching the shell,
answering \cref{rq:patterns}.

% TODO(#2): The systematic search confirmed the SSH gap (\cref{app:protocol}),
% so the SSH part remains an aspect analysis grounded in course observations
% rather than a literature-backed contribution. Revisit whether to keep it as
% a separate contribution or fold it into the others.
